% Generated by application/scripts/105_prepare_joint_qdesn_article_validation_tables_figures.R.
\begin{table}[!htbp]
\centering
\small
\begin{tabular}{@{}>{\raggedright\arraybackslash}p{0.28\textwidth}>{\raggedright\arraybackslash}p{0.62\textwidth}@{}}
\toprule
Item & Value
\\
\midrule
Synthetic scenarios & 9 bridge and stress data-generating mechanisms \\
Scenario classes & bridge (3); stress (6) \\
Quantile grid & 0.05, 0.10, 0.25, 0.50, 0.75, 0.90, 0.95 \\
Series length & 12000 simulated observations; 2000 warmup observations; 10000 retained observations \\
Analysis split & 500 DESN washout observations, 500 fitting observations, 1000 validation observations \\
Forecast protocol & 297 no-refit rolling origins; horizons up to 30 steps ahead; origins spaced 30 observations apart \\
VB controls & VB max iter 240 with adaptive grid 240,480; RHS tau0 = 1 \\
Reported model rows & Joint QDESN RHS; Independent QDESN RHS; Joint exQDESN RHS \\
Held-out diagnostic row & Independent exQDESN RHS, due to localized single-level exAL instability \\
Scoring and diagnostics & True-path MAE/RMSE, check loss, aCRPS, hit-rate error, interval diagnostics, raw crossings and crossings after monotone rearrangement \\
\bottomrule
\end{tabular}
\caption{Protocol for the joint multi-quantile QDESN synthetic validation. All reported rows use the same data-generating realizations, oracle quantiles, fitting sample, validation origins, and scoring rules.}
\label{tab:joint-qdesn-vb-protocol}
\end{table}
